\begin{table}[htbp]\centering
\caption{Sensitivity of the primary estimates to unmeasured confounding:
Oster's $\delta$ and E-values}
\label{tab:sensitivity_bounds}
\small
\begin{tabular}{lcccccc}
\toprule
 & Controlled & No controls & $\delta$ for & $\beta$ at & E-value & E-value \\
 & $\beta$ & $\beta$ & $\beta = 0$ & $\delta = 1$ & (point) & (CI bound) \\
\midrule
Warmth ($W$) & 0.146 & 0.064 & 1.71 & 0.112 & 1.99 & 1.50 \\
Strictness ($S$) & 0.106 & 0.158 & 1.09 & 0.013 & 1.76 & 1.27 \\
\bottomrule
\end{tabular}
\begin{minipage}{\linewidth}
\smallskip
\footnotesize\textit{Notes:} Primary specification on overall Progress~8
($n = 99$); coefficients per standard deviation of the enacted
scores. Oster's $\delta$ (Oster 2019, via \texttt{psacalc}) is the strength of
selection on unobservables, relative to selection on the full control set, that
would drive the coefficient to zero given $R_{\max} = 1.3 \tilde{R}^2 =
0.81$; $\beta$ at $\delta = 1$ is the coefficient that survives when
unobservables are assumed to matter exactly as much as the controls. Each score
is treated as the single treatment with the other score among the controls. The
E-value converts the coefficient to an approximate risk ratio using the outcome
standard deviation on the estimation sample (0.47) and reports the minimum
association a confounder would need with both the score and the outcome to
explain the estimate away; the CI-bound column applies the same conversion to
the boundary of the 95 per cent confidence interval nearer zero. Both measures
address omitted variables only, not reverse causality or measurement error.
\end{minipage}
\end{table}
